% Auto-generated pattern card specifications
% Generated by generate_pattern_cards.py from vocabulary/*.json
% DO NOT EDIT MANUALLY - changes will be overwritten

The following cards show the \textbf{canonical hashed content} for each pattern.
Only semantic fields are included---these are the exact fields that produce the Merkle root hash.
Metadata fields (\texttt{handle}, \texttt{\_meta}, \texttt{tier}, etc.) are \emph{not} part of the hash.

\Needspace{7\baselineskip}
\subsection{sema:StateLock\#mh:SHA-256:c9c2809bdf97bb6ac17a8ffe428d4aa9a28294b99d64cc8a336a239d26a8afb8}

\begin{lstlisting}
{
  "dependencies": {
    "composes_with": {
      "backoff": "sema:Backoff#mh:SHA-256:9e59718422f5408d7892feccb6d538a21103f9573b122b891ee3c8cf66303e72",
      "cooldown": "sema:Cooldown#mh:SHA-256:c288d88055cb0aae80c4b67ed62afcd3ebcc5de18749fc6efb15e9f8c529b9df"
    },
    "references": {
      "actor": "sema:Actor#mh:SHA-256:86fc0930b97ffb4fe85d4820af8e17a43166b745d181895f576edf39edde6aa8",
      "lock": "sema:Lock#mh:SHA-256:95c2ee952a5301d4b346a5e8693350829d88b3c600048eedcf87bb00c23ee5fb",
      "state": "sema:State#mh:SHA-256:4f25e6a16ccd37dc0bc154cccc7055cb504d49daa8b712f372a425eb6c7b8cb2"
    }
  },
  "signature": [
    "Lock(State)"
  ],
  "mechanism": "A coordination pattern where two {{actor}}s temporarily 'fuse' a subset of their writable {{state}}. During the {{lock}}, changes require a cryptographic signature from both. Contention triggers {{backoff}} and {{cooldown}}.",
  "gloss": "Atomic coordination via temporary state fusion",
  "invariants": [
    "Both {{actor}}s sign every write to the fused {{state}} subset.",
    "Exclusive over that subset: no third party writes it while the {{lock}} holds.",
    "Temporary: the fusion ends on completion or timeout, never implicitly."
  ],
  "failure_modes": [
    "Deadlock: both parties wait for each other indefinitely.",
    "Livelock: lock acquired and released rapidly, denying access to others.",
    "Premature dissolution: state auto-dissolves while legitimate work is in progress."
  ]
}
\end{lstlisting}

\Needspace{7\baselineskip}
\subsection{sema:FractalIntelligence\#mh:SHA-256:1dca30e836df7d12d623a804e0eb2aaa37a37bf5247ac0eb9a0c562f262ce67b}

\begin{lstlisting}
{
  "dependencies": {
    "composes_with": {
      "conceptual_decomposition": "sema:ConceptualDecomposition#mh:SHA-256:d0fc6e67c984608fb949bed69ba7e5e3bc022b45bfe173042eec4846b33a5538",
      "localized_learning": "sema:LocalizedLearning#mh:SHA-256:a6dbe910a57206b1a10d3b60733dfe3b98fa85696e701674abe860ee0b730950",
      "marginal_value_rule": "sema:MarginalValueRule#mh:SHA-256:661109205c507b2c26d683d091de7f8b782bf2d6c2e9192373b242a9b078ac71",
      "polymorphic_solver": "sema:PolymorphicSolver#mh:SHA-256:f7fe4635a7520f918c823f0047d6fed8e38fd820c1970ef4623cc8fd78d2b0bd",
      "problem_framer": "sema:ProblemFramer#mh:SHA-256:ad6b7fc1f14f2e26dd0d230f81382a2c602c365ebe4b9ad5fd390fb0c2d4c4fc",
      "reason": "sema:Reason#mh:SHA-256:81ec8196c7c1c761fe353c8d26463bcb61b54f844aa3422a40019f82ef202321",
      "recursion_dive": "sema:RecursionDive#mh:SHA-256:a86f17bad949a7a5405f04b4bc43a3e6be28c5c31e8328ccf178fa292719d754",
      "reframe": "sema:Reframe#mh:SHA-256:cbfd4fadde729367ef8f32a68951e5852535b7f7c3edc68a502d9187b8c1d54a",
      "state_snapshot": "sema:StateSnapshot#mh:SHA-256:d75752f332fe14ae3abf2253a45ab180034d2becc3a7e19e3e30da3c94b3c636",
      "synthesis": "sema:Synthesis#mh:SHA-256:4f1fecba4183e6a79bfb3bd8f1981a6b6fb9f3dced92a484a9aa466c93ae249b"
    },
    "references": {
      "agent": "sema:Agent#mh:SHA-256:e127459069739fffc96a71d54ef6e8c97ca53471d5349c16aa72a092a33756cd",
      "conservation": "sema:Conservation#mh:SHA-256:0b32d007b27a9b5201c46931bdbabd09d08ff9d7b01cf49469f61709989507c2",
      "experience_sharding": "sema:ExperienceSharding#mh:SHA-256:c8a197339886a6bbec3201ead589ad7fd37bbfcde46d6470df59ded8aaa80153",
      "specialize": "sema:Specialize#mh:SHA-256:799d587f9e409c7f9fbf13f5dc83a2437068df8a0738f9bcfb910ee00c0c20a1",
      "strategy": "sema:Strategy#mh:SHA-256:3dc598d0c6af12389df445c204e17df03950dbc886656b85700e6c87c31b2eef",
      "system": "sema:System#mh:SHA-256:f8eb318a1113a10fead02fbf14ae433767ac1a69c9579cdb584cc622068e90b3",
      "task": "sema:Task#mh:SHA-256:a85010f508f85055170c396c8fa390125c1206888c565ac10651f3969a339f8d",
      "universal_solver_tree": "sema:UniversalSolverTree#mh:SHA-256:f9fba480ef013c98691c0ccb6cbbd58dab5183e60798c550ff0de093dbd30727"
    }
  },
  "signature": [
    "System(Reason)"
  ],
  "mechanism": "Expansion of cognitive capability through {{conceptual_decomposition}}: a concept (problem or task) is broken into contract-bound sub-concepts, each governed by the same five-surface Solver Contract (Manifest, Execute, Consult, Verify, Feedback) that governs the parent. A few {{agent}}s can assign solver roles to themselves and perform lightweight fractal intelligence for a specific problem \u2014 the resulting structure may persist as a reusable pattern that improves through use, or may be torn down at completion; both are legitimate modes. The unified {{system}} of scalable cognition uses {{reason}} to orchestrate fractal expansion within the {{universal_solver_tree}}. A {{problem_framer}} initiates by formulating a high-level {{strategy}} before assigning a {{polymorphic_solver}} to a {{task}}; the solver executes a {{recursion_dive}} to spawn child nodes, each applying {{specialize}} with {{localized_learning}}, while {{experience_sharding}} and {{synthesis}} preserve global coherence. {{state_snapshot}} provides crash recovery for persistent instances. {{marginal_value_rule}} governs recursion depth. On failure, {{reframe}} restructures the tree.",
  "gloss": "Expansion of cognitive capability through recursive decomposition of concepts into contract-bounded sub-concepts",
  "invariants": [
    "Fractal Self-Similarity: The process at the Root is identical to the process at the Leaf.",
    "Bounded Expansion: Recursion is limited by Economic constraints (Marginal Value).",
    "Memory {{conservation}}: Specialization must not result in the loss of global context."
  ]
}
\end{lstlisting}

\Needspace{7\baselineskip}
\subsection{sema:SpectralTune\#mh:SHA-256:c0aab008d2cd2f653e56508356987477f9207b4cf6435f4416105b3f37bd544a}

\begin{lstlisting}
{
  "dependencies": {
    "accepts": {
      "signal": "sema:Signal#mh:SHA-256:6eaff8a06e08e75279670bc09d88175d130815240fdfb3a21252c0364769c1fc"
    }
  },
  "mechanism": "Instead of sending a message and hoping it is understood, the sender transmits a 'tuning {{signal}}'\u2014a sequence of hash-based challenges representing the semantic context (ontology, assumptions, version). The receiver must 'resonate' by proving they hold the matching semantic context.",
  "gloss": "Verifying ontology alignment before data transfer",
  "invariants": [
    "No payload data is processed before Tune_ACK.",
    "On hash mismatch, halt and escalate rather than retry."
  ],
  "preconditions": [
    "Both agents possess valid hash function H"
  ],
  "postconditions": [
    "Channel is established with matching context hashes, OR connection terminated and escalated"
  ],
  "parameters": [
    {
      "name": "context_chunks",
      "type": "PositiveInteger",
      "range": "unspecified",
      "description": "Prompts to hash"
    },
    {
      "name": "hash_algo",
      "type": "String",
      "range": "{BLAKE3, SHA256}",
      "description": "Hash algorithm for semantic context verification"
    }
  ],
  "failure_modes": [
    "Over-strict halt: a benign version skew or encoding difference produces a hash mismatch and halts a channel whose parties would have understood each other. Hash equality admits no partial match, so any difference is total."
  ]
}
\end{lstlisting}

\Needspace{11\baselineskip}
\subsection{sema:SteelmanCheck\#mh:SHA-256:b004eb6d85a184da7f9bb7fc06d011bc869f6eb6ad1061f30cbbbcc4c2b8f6aa}

\textbf{Note}: This pattern was downgraded to \textbf{Tier~2} following adversarial analysis.


\begin{lstlisting}
{
  "dependencies": {
    "composes_with": {
      "check": "sema:Check#mh:SHA-256:6890b221c38531710528cd84bbf3a957e29d9836dfc9c02669fc0a86709bcb00",
      "critique": "sema:Critique#mh:SHA-256:646b71acb10a6b7f0874f3d5c7977c44d253f29f9c1231c974d85c7fba131980",
      "judge": "sema:Judge#mh:SHA-256:704cbb71588d7a3bafa08b71c863005e31ec4400860b9cb8b094e660da9c7c6c"
    },
    "references": {
      "agent": "sema:Agent#mh:SHA-256:e127459069739fffc96a71d54ef6e8c97ca53471d5349c16aa72a092a33756cd",
      "belief": "sema:Belief#mh:SHA-256:09e1c919d9829e30422f3fffee11c1af4af54ac7894d157ff371ed50b03f69b1",
      "cognitive_bias": "sema:CognitiveBias#mh:SHA-256:c8a4abe2e0ce048be16d25f890a50f8c7095e099b5ab2de23aa059942328f387",
      "decision": "sema:Decision#mh:SHA-256:b1407127acd4a39314fc24b418310e4b95e0347e4e22f376151ab765e2dfd2c7",
      "loop": "sema:Loop#mh:SHA-256:37f953c0b21b9b0c72a639c4e1a54354faa71f07905531d271f6bbf40c341603",
      "robustness": "sema:Robustness#mh:SHA-256:5304c7824b76c6f308bb1a33ba6171d41fd061a5b104f21710bbb6c7a7c85b43"
    }
  },
  "signature": [
    "Check(Robustness)",
    "Critique(Belief)",
    "Judge(Critique)"
  ],
  "mechanism": "Before finalizing a {{decision}} or output, the {{agent}} MUST generate a claim-specific counter-argument to its conclusion. It uses {{critique}} to challenge the underlying {{belief}}, {{judge}} to score the counter-argument, and {{check}} to test the {{decision}}'s {{robustness}} against any counter-argument that meets `strength_threshold`. A sub-threshold counter-argument is strengthened up to `iteration_limit`; if no adequate counter-argument is produced, release is blocked. Once an adequate counter-argument exists, a Falsified robustness result requires revision or discard, Unknown blocks release pending review, and only Verified permits finalization. It prevents confirmation {{cognitive_bias}}. For adversarial contexts, see adversarial steelmanning.",
  "gloss": "Qualify the opposition, then release only if the conclusion remains robust",
  "invariants": [
    "Adequate opposition: before release, {{judge}} scores the counter-argument at or above `strength_threshold`.",
    "The counter-argument identifies the specific claim it contradicts.",
    "Release discipline: only a Verified {{check}} of the {{decision}}'s {{robustness}} after adequate opposition permits release; Falsified requires revision or discard, and Unknown requires review."
  ],
  "preconditions": [
    "A {{decision}} or output exists to challenge.",
    "The conclusion exposes a specific claim or rationale that can be contradicted."
  ],
  "postconditions": [
    "Either bounded failure to produce an adequate counter-argument blocks release, or an adequate counter-argument and terminal {{check}} result are recorded; the {{decision}} is finalized only on Verified."
  ],
  "parameters": [
    {
      "name": "iteration_limit",
      "type": "Integer",
      "range": "[1, 5]",
      "description": "Maximum strengthening attempts before failure to produce adequate opposition blocks release"
    },
    {
      "name": "strength_threshold",
      "type": "Float",
      "range": "[0.7, 0.95]",
      "description": "Minimum {{judge}} score for a counter-argument to qualify as a steelman; it is not the verdict on the {{decision}}"
    }
  ],
  "failure_modes": [
    "Paralysis by analysis (stuck in {{critique}} {{loop}}).",
    "Collusion: Proposer and Critic are same entity; susceptible to Strawman Waltz attack where agent generates weak counter-arguments to easily defeat them."
  ]
}
\end{lstlisting}

\Needspace{7\baselineskip}
\subsection{sema:WhyClimb\#mh:SHA-256:ed20580f4eecde83bd4822747358418d4b4265c40e468f80b870cd63d3817c93}

\begin{lstlisting}
{
  "dependencies": {
    "references": {
      "condition": "sema:Condition#mh:SHA-256:b480f3c4da11ba6d86f6d14e0fb8521268f3c4bb38588460f7ff4253ad8d78d5",
      "problem": "sema:Problem#mh:SHA-256:2580c43677ec78ee7ee428d200662c6096e4eab38c2c5a868154979074cc0f08",
      "recursive_root_cause": "sema:RecursiveRootCause#mh:SHA-256:f132a6796ffbbddea29f4fe56ff3f81442c99390999348154b6f4ee146be4ae6",
      "reframe": "sema:Reframe#mh:SHA-256:cbfd4fadde729367ef8f32a68951e5852535b7f7c3edc68a502d9187b8c1d54a",
      "solution": "sema:Solution#mh:SHA-256:8d55004ffa8f61d6e7d12d588ae4baa9723296dd7fa501759b5e6b1b97e7c9ab"
    }
  },
  "signature": [
    "Reframe(Problem)"
  ],
  "mechanism": "A root-cause analysis protocol where the agent iteratively asks 'Why is this a problem?' to {{reframe}} the {{problem}} and ascend the abstraction hierarchy. The goal is to reach the 'Ceiling'\u2014the highest level where the problem is still actionable but the solution space is maximized. Utilizes {{recursive_root_cause}}.",
  "gloss": "Recursive problem abstraction",
  "invariants": [
    "Ascension: Scope(Level N+1) > Scope(Level N)",
    "Stop {{condition}}: Halt when 'Why' leads to a value judgment or physics constraint"
  ],
  "preconditions": [
    "{{problem}} is well-formed"
  ],
  "postconditions": [
    "{{solution}} space expanded"
  ],
  "failure_modes": [
    "Climbing too high (solving 'Entropy' instead of 'Fix Bug')."
  ]
}
\end{lstlisting}

\Needspace{11\baselineskip}
\subsection{sema:OptimisticSolver\#mh:SHA-256:63aaa2e083c615a967019ce3aa709f431ff624aafce8b90ed6229144c102fda8}

\textbf{Note}: This pattern was downgraded to \textbf{Tier~2} following adversarial analysis.


\begin{lstlisting}
{
  "dependencies": {
    "composes_with": {
      "atomic_bid": "sema:AtomicBid#mh:SHA-256:5ac637b4b3d74bfaea558ee36b1ee6ffcaf46fe26ac988b524e46faf9821535e",
      "compensate": "sema:Compensate#mh:SHA-256:4b87623f38564fcc95b87d2cbf5df2f5ff41ef27e617ccafd76c7b7896b144e7",
      "compute_budget": "sema:ComputeBudget#mh:SHA-256:4091ac7c6677342920dc53808d02830e35c5849f2af1bb9ba7749c576d37478e",
      "pathway_memory": "sema:PathwayMemory#mh:SHA-256:ce2e6ce822404fae0e60a43fc2c6fbe2e4d51ab9aefee1e7a9e21da75893e722",
      "reflexion": "sema:Reflexion#mh:SHA-256:81eb9140703ea95f53556741df731af84c1cada66cc1e0f3a1d88e4de5fe54d8"
    },
    "references": {
      "parallel": "sema:Parallel#mh:SHA-256:d58d6696be11be09651e71307f0ea6f42bb45fac137ed985a61e475144f687a5",
      "polymorphic_solver": "sema:PolymorphicSolver#mh:SHA-256:f7fe4635a7520f918c823f0047d6fed8e38fd820c1970ef4623cc8fd78d2b0bd",
      "rigorous_solver": "sema:RigorousSolver#mh:SHA-256:2203568438553dc7e6331a7336e2bf488fc99de0a39c601ca3f63ba3e8849038"
    }
  },
  "mechanism": "A high-velocity implementation of {{polymorphic_solver}} designed for efficient multi-agent coordination. Requires a {{parallel}} runtime (Actor Model with Mailboxes) to prevent serial deadlock. It explicitly couples the standard Solver lifecycle (Reason -> Solution) with the {{atomic_bid}} protocol. Unlike the base abstraction, this pattern MANDATES that the agent plan and execute in a single turn. It relies on {{reflexion}} and {{compensate}} for error correction rather than pre-action permission. Use {{compute_budget}} to bound resource consumption. Contrast with {{rigorous_solver}} which prioritizes safety over speed. A {{pathway_memory}} accumulates across runs so the optimistic route-selection converges on strategies that historically succeeded under similar conditions.",
  "gloss": "High-velocity solver requiring parallel runtime",
  "invariants": [
    "Turn Atomicity: Must output [BID] and [TOOL] in the same response.",
    "Non-Blocking: Cannot wait for Orchestrator approval on standard tasks."
  ],
  "preconditions": [
    "Runtime supports Asynchronous/Parallel message delivery (Actor Model).",
    "Message Bus is non-blocking (Mailbox pattern)."
  ],
  "failure_modes": [
    "Over-Eager Execution: Solving the wrong problem efficiently because feedback was skipped.",
    "Serial Deadlock: If deployed on a single-threaded (Talking Stick) engine, multiple atomic outputs will be dropped, halting the swarm."
  ],
  "extends": "sema:PolymorphicSolver#mh:SHA-256:f7fe4635a7520f918c823f0047d6fed8e38fd820c1970ef4623cc8fd78d2b0bd"
}
\end{lstlisting}
